% GDZ PPN236018647,native580lower–581upper;HKUST581–582.
\NumberedItem{112}
\indent Calculum Numerico-Exponentialem omnino novum invenimus.
\par\noindent\hfill[1800] Iun. 12.

\medskip\begingroup\small
\indent Es handelt sich vermutlich um die numerischen Berechnungen von Potenzen der Zahl \(e\), die in einem besondern Päckchen gesammelt im Nachlaß vorhanden (Fh, Nr. 2, Kapsel 50) und Werke III, S. 426—431 unter der Überschrift \textit{Sammlung von Rechnungen, vornehmlich solchen, bei denen von meinen Methoden, die}
\par\endgroup

\clearpage\SourceLeaf{581}\ArithmeticHead{552}{L}{TAGEBUCH MIT ERLÄUTERUNGEN.}
\begingroup\small
\renewcommand{\thefootnote}{*)}
\noindent\textit{Factoren grosser Zahlen zu finden, und von den Wolframschen Logarithmentafeln Gebrauch gemacht ist} abgedruckt sind. Es werden daselbst die gesuchten Exponentialgrößen durch Produkte ganzer Zahlen approximiert, und da man die Logarithmen nach \textsc{Wolfram} sehr genau kennt, so bleibt nur noch \(e^\delta\) zu berechnen, wo \(\delta\) eine sehr kleine, sehr genau bekannte Zahl ist.
\par\indent Die \textsc{Wolfram}schen Tafeln\footnote{\textit{Natürliche oder hyperbolische Logarithmen bis auf 48 Decimalstellen, von Herrn Wolfram berechnet.}}, auf die \textsc{Gauss} in der Überschrift ausdrücklich hinweist, sind zum ersten Male in der \textsc{Schulze}schen \textit{Sammlung logarithmischer u.s.w. Tafeln} veröffentlicht worden; einen Abdruck dieser Tafeln hat \textsc{Gauss} im Jahre 1791, als er zum ersten Male in Braunschweig bei Hofe vorgestellt wurde, von dem damaligen braunschweigischen Staatsminister Geheimen Rat \textsc{Feronce v. Rothenkreuz} zum Geschenk erhalten\renewcommand{\thefootnote}{**)}\footnote{Vergl. \textsc{L. Hänselmann}, \textit{K. F. Gauss}, Leipzig 1878, S. 25, 26.}. Dieser Abdruck mit vielen handschriftlichen Eintragungen von \textsc{Gauss} (vergl. oben S. 11, in der Überschrift) befindet sich im \textsc{Gauss}archiv.
\par\noindent\hfill\textsc{Klein.\quad Schlesinger.}
\par\endgroup
